\documentclass[11pt,a4paper]{article}
\usepackage[margin=2.5cm]{geometry}
\usepackage{booktabs}\usepackage{longtable}\usepackage{xcolor}
\usepackage[colorlinks=true,linkcolor=black,urlcolor=blue]{hyperref}
\usepackage{titlesec}\usepackage{enumitem}\usepackage{tcolorbox}
\titleformat{\section}{\Large\bfseries}{\thesection.}{0.6em}{}
\titlespacing{\section}{0pt}{1.3em}{0.5em}
\setlength{\parskip}{0.55em}\setlength{\parindent}{0pt}
\newcommand{\good}[1]{\textcolor{teal!70!black}{\textbf{#1}}}
\newcommand{\bad}[1]{\textcolor{red!65!black}{\textbf{#1}}}
\newtcolorbox{keybox}{colback=gray!7,colframe=gray!45,boxrule=0.5pt,left=8pt,right=8pt,top=6pt,bottom=6pt}
\title{\textbf{Milestone: the network is not the answer}\\[0.35em]
\large What the research found, what I had wrong, and the plan that replaces it}
\author{Alexander Urvancev \and Yan Beletskiy}
\date{Wednesday 9 September 2026, 22:30 CEST}
\begin{document}\maketitle

\begin{keybox}
\textbf{Verdict: do not rebuild the network. Probability a second attempt beats our hand-written
evaluation: about 7\%.} In its place there are \textbf{six changes with measured Elo from
comparable engines}, together worth roughly 100 Elo, four of them nearly free. And a hardware
finding nobody had costed: \textbf{our machine's memory ceiling has been silently deciding which
improvements we were able to test at all.}
\end{keybox}

\section{Three things I told you that were wrong}

\paragraph{1. ``The architecture cannot see kings, and no training fixes that.''}
I said our 768$\rightarrow$128 network was structurally crippled because its inputs carry no king
information. \textbf{This is the configuration other engines beat their hand-written evaluations
with.} op12no2's Cwtch: 768$\rightarrow$128, on par with a 2700-rated hand-crafted evaluation.
Prokopakop: 768$\rightarrow$128, \good{+200 Elo}. Leorik measured a width curve where \textbf{128
is the peak} and 512 is \emph{worse}. The author of the standard training tool writes: ``just
start with basic 768 inputs.''

Our architecture was fine. I diagnosed a fatal flaw that is not one.

\paragraph{2. ``We have an incremental accumulator.''}
We do not. \texttt{fasteval.py:655} rebuilds the whole thing from the piece list on every single
call --- and says so in its own docstring, deliberately, for correctness. I repeated the opposite
several times this evening. It also means the 14.6\% speed cost we measured is a \emph{ceiling},
not a floor.

\paragraph{3. Mobility as the leading candidate for Thursday.}
\bad{Ruled out.} A mobility term is worth $+41 \pm 14$ elsewhere --- but at our measured
1.7$\times$ time-to-depth it costs roughly 60--65 Elo of lost search. \textbf{Net loss.} Other
engines pay almost nothing for it because they count attacks the move generator already
computes; ours would need 4--6 hours of optimisation before the term is even affordable.

\section{Why the network still fails}

\subsection{The one real deviation}

\begin{center}
\begin{tabular}{lr}
\toprule
Positions per parameter, working range across the field & 50--3000 \\
\textbf{Ours} & \bad{1.9} \\
\midrule
Lowest volume anyone reports a working net at & 1.5 M positions \\
\textbf{Ours} & \bad{186{,}190} \\
\bottomrule
\end{tabular}
\end{center}

A master's thesis abandoned \emph{exactly} our approach --- online games labelled by Stockfish ---
with the note that such datasets are ``too small for this task (order of hundreds of millions).''

\subsection{But that is not what cost us 67 Elo}

A starved network is a network that judges positions badly. \textbf{Ours judged them better than
our hand-written evaluation on four independent measurements, including at real search leaves,
and still played 67 Elo worse.} No source anywhere connects a data shortage to that reversal.

And the three fixes the research proposed have each already been tested here and killed: the
outputs are not compressed (they are 5--19\% \emph{wider}), the error does not explode
off-distribution (the network stays ahead at leaves), and the improved loss function targets a
metric \textbf{the network already wins.}

\begin{keybox}
So the honest position is not ``we know what went wrong and can fix it.'' It is
\textbf{``the network was better by every measure we own and played worse, and nobody has
explained why.''} The right response to an unexplained result with 36 hours left is not to
rebuild it ten times larger.
\end{keybox}

\subsection{The arithmetic}

$0.7$ (pipeline correct first time, no iteration) $\times$ $0.35$ (true effect positive at all)
$\times$ $0.3$ (a screen we can afford resolves it) $\approx$ \bad{7\%}.

And a warning worth heeding: our promotion rule has already slipped twice on rested judgement
with nothing sunk. At 04:00 Friday after thirty hours on a network, it slips again --- and it
slips toward shipping.

\section{The hardware finding}

\begin{center}
\begin{tabular}{lr}
\toprule
Our machine & 16 threads, 7 GB RAM, \bad{1 GB available} \\
Transposition table (\texttt{TT\_BITS = 21}) & 2{,}097{,}152 slots $=$ 59 MB \\
Nodes pushed through it in one move & \bad{1.6--3.2 million} \\
Platform allows per process & 2 GB \\
\bottomrule
\end{tabular}
\end{center}

\textbf{Our memory of past positions is smaller than a single move's search.} Raising it one step
is a one-line change --- but 235 MB $\times$ 24 agent processes is 5.6 GB, on a box with 1 GB
free. \textbf{It has never been tested because our machine physically cannot test it.}

That is the real reason to rent a box, and it is a better one than the network: \textbf{a second
machine roughly doubles the number of measured decisions we can make before Friday}, and the
number of measured decisions is the binding constraint on this project.

\section{The plan that replaces it}

Every figure below is measured in a comparable engine, with its uncertainty.

\begin{center}
\begin{longtable}{p{4.6cm}rrp{5.4cm}}
\toprule
Change & Elo & Hours & Note \\
\midrule\endhead
\textbf{SEE} (already screening) & $+26$ & \good{0} & Running now, result 06:20. Makes the search
\emph{faster}, and targets our blunder rate. \\
\addlinespace
\textbf{Internal iterative reduction} & $+9.7 \pm 5.5$ & 0.3 & Three lines. Best Elo per line
available. \\
\addlinespace
\textbf{Transposition table} $21\to23$ & unpublished & 0.1 & Direction certain. Needs the rented
box to test at all. \\
\addlinespace
\textbf{History gravity + malus} & $+15.9 \pm 7.4$ & 1 & $\sim$15 lines, no measurable cost. \\
\addlinespace
\textbf{Correction history} & $+18.9 \pm 8.0$ & 3 & \textbf{The one that targets what we
measured.} \\
\addlinespace
\textbf{Time management} & $+28$ to $+43$ & 2 & Real, but two constants must move together or we
risk losing on time. \\
\bottomrule
\end{longtable}
\end{center}

\subsection{Why correction history is the pick}

We measured that at real search leaves our evaluation is off by a median of \textbf{304--342
centipawns} --- about a minor piece --- while our pruning margins are 150--750. We are making
decisions on differences smaller than our own error.

Correction history records, for each material configuration, \textbf{how far the search's verdict
usually lands from the static evaluation}, and adds that learned offset back before the pruning
decisions are made. It is the only item on the list that attacks the pathology we actually
observed rather than one we assumed.

It costs one array lookup. It passed 128{,}672 games of testing in Stockfish and \emph{scales up}
with longer time controls --- and ours is the long one.

\section{Timetable}

\begin{center}
\begin{tabular}{lll}
\toprule
Thu 06:20 & SEE result lands & decides the base build \\
Thu 06:20--10:00 & Code: table size, iterative reduction, history malus & 1.4 h \\
Thu, in parallel & Correction history & 3 h \\
Thu 10:00--22:00 & Two screens per machine, two machines & $\sim$12 h \\
Thu 22:00 & \textbf{Freeze}, 500-game safety gate & 1 h \\
Fri 08:00 & \textbf{Upload} (early submission is a tie-break) & \\
Fri 11:00 & Locked & \\
\bottomrule
\end{tabular}
\end{center}

Three hours of margin before close. The network plan fitted the window with zero slack, on a
throughput figure nobody had measured, and consumed the entire second machine for its whole life.

\section{What this changes about winning}

Six measured changes worth perhaps 100 Elo beats one unexplained gamble at 7\%. That is the whole
argument.

It also means the second machine is worth buying \emph{more}, not less --- just for a different
job than the one we were discussing an hour ago. Its work is to double how many of those six we
can actually prove before Friday.

\end{document}
